% Image credits for the photographs and AI illustrations of Book 3
% (University Chemistry, Year 2). Machine-readable license records live in
% images/book3/CREDITS.md; keep the two files in sync. The Book 3 writer
% owns this file: add one \item per photograph (in an itemize, after the
% first paragraph) as photographs land.
\chapter*{Créditos das imagens}

Os desenhos deste livro são originais. As fotografias, quando usadas, são
utilizadas sob as respectivas licenças livres, com agradecimento aos seus
autores; os links completos para as fontes e as URLs das licenças de cada
fotografia estão listados em \texttt{images/book3/CREDITS.md} no
repositório do livro.
Os pictogramas de perigo são os do Sistema Globalmente Harmonizado de
Classificação e Rotulagem de Produtos Químicos, tal como desenhados para o
pacote \texttt{ghsystem} do \TeX\ Live.

\bigskip
Ao lado das fotografias e dos desenhos originais, algumas figuras são
ilustrações geradas por IA, produzidas com a geração de imagens da OpenAI a
partir de instruções escritas pelos editores do livro, e revisadas quanto à
exatidão química antes de serem incluídas. A instrução completa de cada
imagem gerada está registrada em \texttt{images/book3/ai/PROMPTS.md} no
repositório.

\bigskip
\noindent Fotografias, por capítulo:
\begin{itemize}\small
  \item Cap.~1: Germain Henri Hess, retrato por I.~I. Reimers (1837--1838),
        CC BY-SA 4.0; Marcellin Berthelot, autor desconhecido, domínio público
        (ambas da Wikimedia Commons).
  \item Cap.~2: Josiah Willard Gibbs, fotógrafo desconhecido, domínio público
        (Wikimedia Commons).
  \item Cap.~3: François-Marie Raoult, fotógrafo desconhecido, domínio público
        (Wikimedia Commons).
  \item Cap.~4: Jacobus Henricus van 't Hoff, por Nicola Perscheid (1904),
        domínio público; Henry Le Chatelier, autor desconhecido, CC BY-SA 4.0;
        Fritz Haber e Carl Bosch, Fundação Nobel, domínio público (todas da
        Wikimedia Commons).
  \item Cap.~6: alto-forno de Duisburg-Nord, de Dietmar Rabich, CC BY-SA
        4.0; soldagem aluminotérmica de um trilho, de PetrS., CC BY-SA 3.0
        (ambas da Wikimedia Commons).
  \item Cap.~7: coluna de destilação de uma refinaria, de CEphoto, Uwe Aranas,
        CC BY-SA 3.0 (Wikimedia Commons).
  \item Cap.~9: Michael Faraday, por Thomas Phillips (1842), domínio público
        (Wikimedia Commons).
  \item Cap.~10: Jaroslav Heyrovsk\'y, arquivo do Instituto J.~Heyrovsk\'y
        de Físico-Química, CC BY-SA 3.0 (Wikimedia Commons).
  \item Cap.~11: uma pilha de Volta, de GuidoB, CC BY-SA 3.0; Charles Martin Hall
        e Paul H\'eroult, autores desconhecidos, domínio público (todas da
        Wikimedia Commons).
  \item Cap.~12: uma estátua de cobre, de Elcobbola, domínio público; um ânodo
        de sacrifício, de Zwergelstern, CC BY-SA 3.0 (ambas da Wikimedia Commons).
  \item Cap.~13: Erwin Schr\"odinger, Fundação Nobel (1933), domínio público
        (Wikimedia Commons).
  \item Cap.~14: oxigênio líquido em um ímã, de Pieter Kuiper, domínio público;
        Robert Mulliken e Friedrich Hund, de GFHund, CC BY 3.0 (ambas da
        Wikimedia Commons).
  \item Cap.~16: August Kekul\'e (1873), autor desconhecido, domínio público
        (Wikimedia Commons).
  \item Cap.~18: soluções de metais de transição, de Benjah-bmm27, domínio
        público; Alfred Werner, Les Prix Nobel (1914), domínio público (ambas da
        Wikimedia Commons).
  \item Cap.~19: um cristal de rubi, domínio público; uma esmeralda, de Didier
        Descouens, CC BY-SA 3.0; cristais de sulfato de cobre, de Stephanb,
        CC BY-SA 3.0 (todas da Wikimedia Commons).
  \item Cap.~20: Akira Suzuki, Ei-ichi Negishi e Richard Heck, de BloodIce,
        CC BY-SA 4.0 (Wikimedia Commons).
  \item Cap.~22: Charles Friedel (anos 1890) e James Mason Crafts (\textit{Popular
        Science Monthly}, 1900), domínio público (Wikimedia Commons).
  \item Cap.~24: Michel Eug\`ene Chevreul, por Nicolas Eustache Maurin, domínio
        público (Wikimedia Commons).
  \item Cap.~25: Alexander Borodin, autor desconhecido, domínio público; Ludwig
        Claisen, de Veronica.villagrasa, CC BY-SA 4.0 (ambas da Wikimedia Commons).
  \item Cap.~26: Robert Robinson, autor desconhecido, domínio público
        (Wikimedia Commons).
  \item Cap.~28: Elias James Corey, de Trvthchem, CC0 (Wikimedia Commons).
  \item Cap.~29: Wallace Carothers no laboratório, fotógrafo desconhecido,
        domínio público (Wikimedia Commons).
  \item Cap.~30: Emil Fischer, Atelier Victoria, Berlim, por volta de 1895,
        domínio público (Wikimedia Commons).
  \item Cap.~31: Mikhail Tsvet, autor desconhecido, domínio público (Wikimedia
        Commons).
  \item Cap.~32: Francis William Aston (1922), autor desconhecido, domínio
        público; Gustav Kirchhoff, Robert Bunsen e Henry Roscoe (1862), domínio
        público (ambas da Wikimedia Commons).
  \item Cap.~34: William Sealy Gosset (1908), domínio público (Wikimedia
        Commons).
%% next chapter here
\end{itemize}
\clearpage
